\documentclass[11pt,fontset=fandol]{ctexart}
\usepackage[a4paper,margin=23mm]{geometry}
\usepackage{amsmath,amssymb,mathtools,bm,mathrsfs,esint}
\usepackage{graphicx,longtable,booktabs,array,calc}
\usepackage[normalem]{ulem}
\usepackage{xcolor,hyperref}
\hypersetup{hidelinks}
\setcounter{secnumdepth}{-1}
\newcounter{none}
\providecommand{\tightlist}{\setlength{\itemsep}{0pt}\setlength{\parskip}{0pt}}
\providecommand{\passthrough}[1]{#1}
\usepackage{amssymb,mathrsfs,esint,bm,newunicodechar}
\allowdisplaybreaks
\setlength{\emergencystretch}{3em}
\providecommand{\llangle}{\langle\!\langle}
\providecommand{\rrangle}{\rangle\!\rangle}
\newunicodechar{∞}{\ensuremath{\infty}}
\newunicodechar{∈}{\ensuremath{\in}}
\newunicodechar{∑}{\ensuremath{\sum}}
\newunicodechar{∫}{\ensuremath{\int}}
\newunicodechar{≤}{\ensuremath{\leq}}
\newunicodechar{≥}{\ensuremath{\geq}}
\newunicodechar{≠}{\ensuremath{\neq}}
\newunicodechar{→}{\ensuremath{\to}}
\newunicodechar{⇒}{\ensuremath{\Rightarrow}}
\newunicodechar{α}{\ensuremath{\alpha}}
\newunicodechar{β}{\ensuremath{\beta}}
\newunicodechar{γ}{\ensuremath{\gamma}}
\newunicodechar{λ}{\ensuremath{\lambda}}
\newunicodechar{φ}{\ensuremath{\varphi}}
\newunicodechar{ε}{\ensuremath{\varepsilon}}
\newunicodechar{σ}{\ensuremath{\sigma}}
\newunicodechar{θ}{\ensuremath{\theta}}
\newunicodechar{δ}{\ensuremath{\delta}}
\newunicodechar{η}{\ensuremath{\eta}}
\newunicodechar{π}{\ensuremath{\pi}}
\newunicodechar{ξ}{\ensuremath{\xi}}
\newunicodechar{Σ}{\ensuremath{\Sigma}}
\newunicodechar{Δ}{\ensuremath{\Delta}}
\newunicodechar{∂}{\ensuremath{\partial}}
\newunicodechar{∀}{\ensuremath{\forall}}
\newunicodechar{∃}{\ensuremath{\exists}}
\newunicodechar{√}{\ensuremath{\surd}}
\newunicodechar{−}{\ensuremath{-}}
\newunicodechar{μ}{\ensuremath{\mu}}
\newunicodechar{ρ}{\ensuremath{\rho}}
\newunicodechar{τ}{\ensuremath{\tau}}
\newunicodechar{ω}{\ensuremath{\omega}}
\newunicodechar{Ω}{\ensuremath{\Omega}}
\newunicodechar{χ}{\ensuremath{\chi}}
\newunicodechar{ψ}{\ensuremath{\psi}}
\newunicodechar{Γ}{\ensuremath{\Gamma}}
\newunicodechar{Φ}{\ensuremath{\Phi}}
\newunicodechar{Λ}{\ensuremath{\Lambda}}
\newunicodechar{ℝ}{\ensuremath{\mathbb{R}}}
\newunicodechar{ℤ}{\ensuremath{\mathbb{Z}}}
\newunicodechar{ℕ}{\ensuremath{\mathbb{N}}}
\newunicodechar{△}{\ensuremath{\triangle}}
\newunicodechar{⊗}{\ensuremath{\otimes}}
\begin{document}
\section{2013年杭州电子科技大学601数学分析考研真题参考答案}\label{ux5e74ux676dux5ddeux7535ux5b50ux79d1ux6280ux5927ux5b66601ux6570ux5b66ux5206ux6790ux8003ux7814ux771fux9898ux53c2ux8003ux7b54ux6848}

来源：\texttt{HDU/601数学分析/2013/2013年杭州电子科技大学601数学分析考研真题参考答案.pdf}。原
PDF 页序：1, 2, 3, 4, 5, 6, 7, 8。

\section{杭州电子科技大学2013年数学分析真题解析}\label{ux676dux5ddeux7535ux5b50ux79d1ux6280ux5927ux5b662013ux5e74ux6570ux5b66ux5206ux6790ux771fux9898ux89e3ux6790}

\begin{quote}
助手整理：依据同题真题、原答案中可相互核对的公式及直接计算恢复数学内容。推导经重新核算，不作为原参考答案的逐字转录；题面与答案存在分歧时明确说明。
\end{quote}

\subsection{一、填空题}\label{ux4e00ux586bux7a7aux9898}

\subsubsection{1. 指数型极限}\label{ux6307ux6570ux578bux6781ux9650}

\[
\frac{2x}{1-x}\ln\frac{2x}{1+x}
\sim\frac{2x}{1-x}\left(\frac{2x}{1+x}-1\right)
=-\frac{2x}{1+x}\longrightarrow-1.
\] 所以原极限为 \(\boxed{e^{-1}}\)。

\subsubsection{2. 不定积分}\label{ux4e0dux5b9aux79efux5206}

令 \(t=1/x\)，则 \[
\boxed{\int\sin(1/x)\frac{dx}{x^2}=\cos(1/x)+C.}
\]

\subsubsection{3. 全微分}\label{ux5168ux5faeux5206}

答案中的连续求导公式确认函数为 \(z=e^{\sin(xy)}\)，故 \[
\boxed{dz=e^{\sin(xy)}\cos(xy)(y\,dx+x\,dy).}
\]

\subsubsection{4. 换序积分}\label{ux6362ux5e8fux79efux5206}

题目及答案区域对应 \(0\leq x\leq1\)、\(x^2\leq y\leq x\)，换序后为
\(0\leq y\leq1\)、\(y\leq x\leq\sqrt y\)。因此 \[
\int_0^1dx\int_{x^2}^xf(x,y)\,dy
=\int_0^1dy\int_y^{\sqrt y}f(x,y)\,dx,\qquad
\boxed{\varphi(y)=\sqrt y.}
\]

\subsubsection{5. 微分}\label{ux5faeux5206}

对 \(y=(1+\sin x)^x\)， \[
dy=(1+\sin x)^x\left[\ln(1+\sin x)+\frac{x\cos x}{1+\sin x}\right]dx.
\] 所以 \(\boxed{\left.dy\right|_{x=\pi}=-\pi\,dx}\)。

\subsubsection{6.
分段函数连续}\label{ux5206ux6bb5ux51fdux6570ux8fdeux7eed}

同题答案的左右极限确认函数为 \[
f(x)=\begin{cases}\sin(ax)/x,&x<0,\\2-e^x,&x\geq0.\end{cases}
\] 两侧极限分别为 \(a\)、1，且 \(f(0)=1\)，故 \(\boxed{a=1}\)。

\subsubsection{7.
积分关系求导}\label{ux79efux5206ux5173ux7cfbux6c42ux5bfc}

对 \(\int_1^{x^2}f(t)\,dt=\int_1^x\varphi(t)\,dt\) 两边求导，得到 \[
\boxed{\varphi(x)=2xf(x^2).}
\]

\subsubsection{8. 级数求和}\label{ux7ea7ux6570ux6c42ux548c}

由 \(\sum_{n\geq0}x^n=(1-x)^{-1}\)，在 \(|x|<1\) 内求导再乘 \(x\)，得 \[
\sum_{n=1}^{\infty}nx^n=\frac{x}{(1-x)^2}.
\] 取 \(x=1/2\)，得所求和为 \(\boxed2\)。

\subsection{二、计算题}\label{ux4e8cux8ba1ux7b97ux9898}

\subsubsection{1. 极限}\label{ux6781ux9650}

\[
\frac1{\ln(1+x)}-\frac1x
=\frac{x-\ln(1+x)}{x\ln(1+x)}\longrightarrow\boxed{\frac12}.
\] 这里使用 \(x-\ln(1+x)\sim x^2/2\)。

\subsubsection{2. 定积分}\label{ux5b9aux79efux5206}

令 \(t=\sqrt{e^x-1}\)，则 \(x=\ln(1+t^2)\)，上下限由 \(0,\ln2\) 变为
\(0,1\)。因此 \[
\boxed{\int_0^{\ln2}\sqrt{e^x-1}\,dx
=\int_0^1\frac{2t^2}{1+t^2}\,dt
=2[t-\arctan t]_0^1=2-\frac\pi2.}
\]

\subsubsection{3.
二阶偏导数组合}\label{ux4e8cux9636ux504fux5bfcux6570ux7ec4ux5408}

原真题要求计算的是 \(xu_{xx}+yu_{xy}\)，其中 \[
u=yf(x/y)+xg(y/x).
\] 部分答案 OCR 将目标误识为
\(xu_{xx}+yu_{yy}\)。这里严格采用原真题表达式。在 \(xy\neq0\)
的定义域内， \[
u_{xx}=\frac1y f''(x/y)+\frac{y^2}{x^3}g''(y/x),
\] \[
u_{xy}=-\frac{x}{y^2}f''(x/y)-\frac{y}{x^2}g''(y/x).
\] 所以 \[
\boxed{xu_{xx}+yu_{xy}=0.}
\] 也可由一次齐次函数的 Euler 恒等式 \(xu_x+yu_y=u\) 对 \(x\) 求导得到。

\subsubsection{4.
圆锥与圆柱围成的体积}\label{ux5706ux9525ux4e0eux5706ux67f1ux56f4ux6210ux7684ux4f53ux79ef}

圆柱内部的极坐标范围为
\(0\leq\theta\leq\pi\)、\(0\leq r\leq2\sin\theta\)，圆锥之间的竖直范围为
\(-|a|r\leq z\leq|a|r\)。因此 \[
\begin{aligned}
V&=2|a|\int_0^\pi\int_0^{2\sin\theta}r^2\,dr\,d\theta\\
&=\frac{16|a|}{3}\int_0^\pi\sin^3\theta\,d\theta
=\boxed{\frac{64|a|}{9}}.
\end{aligned}
\] 当 \(a>0\) 时即为 \(64a/9\)。答案中只算上半部分得到的 \(32a/9\)
还需要乘 2；原题没有给 \(a>0\) 时应保留绝对值。

\subsection{三、综合题}\label{ux4e09ux7efcux5408ux9898}

\subsubsection{\texorpdfstring{1. \(f(x)=\ln(1+x^2)\)
的性质}{1. f(x)=\textbackslash ln(1+x\^{}2) 的性质}}\label{fxln1x2-ux7684ux6027ux8d28}

\[
f'(x)=\frac{2x}{1+x^2},\qquad f''(x)=\frac{2(1-x^2)}{(1+x^2)^2}.
\] 函数在 \((-\infty,0]\) 上严格递减，在 \([0,+\infty)\)
上严格递增，\(f(0)=0\) 是极小值，也是全局最小值，无极大值。

当 \(|x|<1\) 时 \(f''>0\)，当 \(|x|>1\) 时 \(f''<0\)，两个拐点为
\(\boxed{(-1,\ln2),(1,\ln2)}\)。

\subsubsection{2.
级数的敛散性}\label{ux7ea7ux6570ux7684ux655bux6563ux6027}

对 \(x\neq0\)，通项绝对值的比值极限为 \[
\lim_{n\to\infty}\frac{\sin(\pi/2^{n+1})}{|x|\sin(\pi/2^n)}=\frac1{2|x|}.
\] 因此 \(|x|>1/2\) 时绝对收敛，\(0<|x|<1/2\) 时发散。在 \(x=\pm1/2\)
时，通项绝对值趋于 \(\pi\)，故也发散；\(x=0\) 时原式无定义。收敛域为 \[
\boxed{(-\infty,-1/2)\cup(1/2,+\infty).}
\]

\subsubsection{3.
上半圆曲线积分}\label{ux4e0aux534aux5706ux66f2ux7ebfux79efux5206}

本题与 2011 年二（4）的被积形式、路径相同： \[
I=\int_{C^+}d(e^{-x^2}\sin y)+\int_{C^+}x^4\,dy.
\] 端点函数值均为零；令 \(x=\cos t,y=\sin t\)、\(0\leq t\leq\pi\)，得到
\[
\boxed{I=\int_0^\pi\cos^5t\,dt=0.}
\]

\subsection{四、证明题}\label{ux56dbux8bc1ux660eux9898}

\subsubsection{1. 对数不等式}\label{ux5bf9ux6570ux4e0dux7b49ux5f0f}

对 \(x>0\)， \[
0<x-\ln(1+x)=\int_0^x\frac{t}{1+t}\,dt<\int_0^xt\,dt=\frac{x^2}{2}.
\] 移项得到 \(\boxed{x-x^2/2<\ln(1+x)<x}\)。

\subsubsection{2.
一致连续函数列的一致极限}\label{ux4e00ux81f4ux8fdeux7eedux51fdux6570ux5217ux7684ux4e00ux81f4ux6781ux9650}

给定 \(\varepsilon>0\)，由一致收敛可取 \(N\) 使
\(\sup_{x\in I}|f_N(x)-f(x)|<\varepsilon/3\)。由 \(f_N\) 一致连续，存在
\(\delta>0\)，使 \(x,y\in I\) 且 \(|x-y|<\delta\) 时有
\(|f_N(x)-f_N(y)|<\varepsilon/3\)。于是 \[
|f(x)-f(y)|\leq|f(x)-f_N(x)|+|f_N(x)-f_N(y)|+|f_N(y)-f(y)|<\varepsilon.
\] 故 \(f\) 在 \(I\) 上一致连续。

\subsubsection{3.
原点处连续但不可微}\label{ux539fux70b9ux5904ux8fdeux7eedux4f46ux4e0dux53efux5fae}

\[
f(x,y)=\begin{cases}\dfrac{x^2y^2}{(x^2+y^2)^{3/2}},&(x,y)\neq(0,0),\\0,&(x,y)=(0,0).\end{cases}
\] 因为 \(x^2y^2\leq(x^2+y^2)^2/4\)， \[
0\leq f(x,y)\leq\frac14\sqrt{x^2+y^2}\longrightarrow0,
\] 所以 \(f\) 在原点连续。由坐标轴上的函数值为零可得
\(f_x(0,0)=f_y(0,0)=0\)。若可微，则应有
\(f(x,y)/\sqrt{x^2+y^2}\to0\)；但沿 \(y=x\neq0\)， \[
\frac{f(x,x)}{\sqrt{2x^2}}=\frac14.
\] 故 \(f\) 在原点不可微。

\subsubsection{4. Gamma
型积分收敛}\label{gamma-ux578bux79efux5206ux6536ux655b}

当 \(p>0\) 时，\(x\to0^+\) 有
\(x^{p-1}e^{-x}\sim x^{p-1}\)，因此在零点附近可积。又因
\(x^{p-1}e^{-x/2}\to0\)，存在 \(M\)，使 \(x\geq M\) 时 \[
0\leq x^{p-1}e^{-x}\leq e^{-x/2}.
\] 右端在无穷远处可积，中间有限区间上被积函数连续，所以
\(\int_0^{\infty}x^{p-1}e^{-x}\,dx\) 收敛。

\subsubsection{5.
中值点的存在性}\label{ux4e2dux503cux70b9ux7684ux5b58ux5728ux6027}

令 \(F(x)=xf(x)\)。由积分中值定理，存在 \(c\in(0,1)\)，使 \[
F(c)=\int_0^1F(x)\,dx=f(1)=F(1).
\] 对 \(F\) 在 \([c,1]\) 上应用 Rolle 定理，存在 \(\xi\in(c,1)\)，满足
\[
0=F'(\xi)=f(\xi)+\xi f'(\xi).
\] 因为 \(\xi>0\)，得到 \(\boxed{f'(\xi)=-f(\xi)/\xi}\)。
\end{document}
